\documentclass[10pt,twoside]{article}
\usepackage[utf8]{inputenc}
\usepackage[vietnamese]{babel}
\usepackage{amsmath,amssymb,amsfonts}
\usepackage{graphicx}
\usepackage{xcolor}
\usepackage{tcolorbox}
\usepackage{geometry}
\usepackage{fancyhdr}
\usepackage{mathpazo}
\usepackage{helvet}

\geometry{
    paperwidth=8.5in,
    paperheight=11in,
    top=0.7in,
    bottom=0.8in,
    left=0.75in,
    right=0.75in
}

\definecolor{stewartcyan}{RGB}{0, 136, 206}
\definecolor{stewartgreen}{RGB}{0, 130, 70}
\definecolor{stewartred}{RGB}{195, 30, 45}
\definecolor{copyrightgray}{RGB}{110, 110, 110}

\pagestyle{fancy}
\fancyhf{}
\renewcommand{\headrulewidth}{0pt}
\fancyhead[RO,LE]{{\small\sffamily\bfseries\color{gray!80!black}MỤC 3.5}\hspace{1.2em}{\small\sffamily Đạo hàm ẩn}\hspace{2.5em}{\normalsize\bfseries\sffamily 213}}

\fancyfoot[C]{%
    \fontsize{6.5pt}{8pt}\selectfont\color{copyrightgray}%
    \textcopyright\,2021 Cengage Learning. Bản quyền đã được bảo lưu. Không được sao chép, quét hoặc nhân bản toàn bộ hay một phần. WCN 02-200-203\\
    Đánh giá của biên tập viên cho rằng nội dung bị lược bỏ không ảnh hưởng đáng kể đến trải nghiệm học tập tổng thể. Cengage Learning bảo lưu quyền xóa nội dung bổ sung vào bất kỳ lúc nào nếu có yêu cầu hạn chế về bản quyền sau đó.%
}

\newtcolorbox{historybox}{
    colback=white,
    colframe=stewartgreen,
    arc=0mm,
    boxrule=0.75pt,
    left=6pt,
    right=6pt,
    top=6pt,
    bottom=6pt
}

\setlength{\parindent}{0pt}
\setlength{\parskip}{4pt}

\begin{document}

\noindent
\begin{minipage}[t]{0.27\textwidth}
    \vspace{0pt}
    \begin{historybox}
        {\noindent\bfseries\sffamily\color{stewartgreen}\small Abel và Galois}\par\vspace{3pt}
        {\fontsize{7.8pt}{10.2pt}\selectfont
        Nhà toán học người Na Uy Niels Abel đã chứng minh vào năm 1824 rằng không thể có một công thức tổng quát nào biểu diễn qua căn thức cho các nghiệm của phương trình $p(x) = 0$, trong đó $p$ là một đa thức bậc 5 với các hệ số nguyên.\par\vspace{3pt}
        Sau đó, nhà toán học người Pháp Évariste Galois đã chứng minh rằng không thể tìm được một công thức tổng quát biểu diễn qua căn thức cho các nghiệm của một phương trình $p(x) = 0$, trong đó $p$ là một đa thức bậc $n \ge 5$.\par}
    \end{historybox}

    \vspace{1.6em}
    \begin{center}
        \includegraphics[width=0.92\linewidth]{images/p0248_fig1.png}\\[0.4em]
        {\fontsize{8.5pt}{10.5pt}\selectfont\textbf{\textsf{HÌNH 6}}}
    \end{center}

    \vspace{1.8em}
    \begin{center}
        \includegraphics[width=0.95\linewidth]{images/p0248_fig2.png}\\[0.4em]
        {\fontsize{8.5pt}{10.5pt}\selectfont\textbf{\textsf{HÌNH 7}}\\[0.2em]
        $(x^2 - 1)(x^2 - 4)(x^2 - 9)$\\
        $= y^2(y^2 - 4)(y^2 - 9)$}
    \end{center}
\end{minipage}%
\hfill
\begin{minipage}[t]{0.69\textwidth}
    \vspace{0pt}
    (Đây là ba hàm số có đồ thị được thể hiện trong Hình~3.) Bạn có thể thấy rằng phương pháp đạo hàm ẩn giúp tiết kiệm một khối lượng công việc khổng lồ trong những trường hợp như thế này. Hơn nữa, đạo hàm ẩn cũng áp dụng dễ dàng tương tự cho các phương trình chẳng hạn như
    \begin{equation*}
    y^5 + 3x^2y^2 + 5x^4 = 12
    \end{equation*}
    mà ở đó việc tìm một biểu thức cho $y$ theo $x$ là \textit{bất khả thi}.

    \vspace{0.6em}
    \noindent\textbf{\textsf{\color{stewartcyan}VÍ DỤ 3}}\hspace{1em}Tìm $y'$ nếu $\sin(x + y) = y^2 \cos x$.

    \vspace{0.3em}
    \noindent\textbf{\textsf{\color{stewartcyan}LỜI GIẢI}}\hspace{1em}Lấy đạo hàm ẩn hai vế theo $x$ và ghi nhớ rằng $y$ là một hàm của $x$, ta được
    \begin{equation*}
    \cos(x + y) \cdot (1 + y') = y^2(-\sin x) + (\cos x)(2yy')
    \end{equation*}
    (Lưu ý rằng chúng ta đã sử dụng Quy tắc chuỗi ở vế trái và Quy tắc nhân cùng Quy tắc chuỗi ở vế phải.) Nếu thu gọn các số hạng có chứa $y'$, ta được
    \begin{equation*}
    \cos(x + y) + y^2 \sin x = (2y \cos x)y' - \cos(x + y) \cdot y'
    \end{equation*}
    Do đó
    \begin{equation*}
    y' = \frac{y^2 \sin x + \cos(x + y)}{2y \cos x - \cos(x + y)}
    \end{equation*}
    Hình~6, được vẽ bằng máy tính, thể hiện một phần của đường cong $\sin(x + y) = y^2 \cos x$. Để kiểm tra lại phép tính của chúng ta, chú ý rằng $y' = -1$ khi $x = y = 0$ và từ đồ thị ta thấy độ dốc xấp xỉ bằng $-1$ tại gốc tọa độ.\hfill{\color{stewartcyan}$\blacksquare$}

    \vspace{0.6em}
    Các Hình~7, 8 và 9 thể hiện thêm ba đường cong nữa được vẽ bởi máy tính. Trong Bài tập 45--46, bạn sẽ có cơ hội tạo ra và khảo sát những đường cong kỳ thú có dạng này.

    \vspace{0.8em}
    \noindent
    \begin{minipage}[t]{0.48\linewidth}
        \centering
        \includegraphics[width=0.96\linewidth]{images/p0248_fig3.png}\\[0.4em]
        {\fontsize{8.5pt}{10.5pt}\selectfont\textbf{\textsf{HÌNH 8}}\\[0.2em]
        $\cos(x - \sin y) = \sin(y - \sin x)$}
    \end{minipage}%
    \hfill
    \begin{minipage}[t]{0.48\linewidth}
        \centering
        \includegraphics[width=0.96\linewidth]{images/p0248_fig4.png}\\[0.4em]
        {\fontsize{8.5pt}{10.5pt}\selectfont\textbf{\textsf{HÌNH 9}}\\[0.2em]
        $\sin(xy) = \sin x + \sin y$}
    \end{minipage}

    \vspace{1.2em}
    \noindent{\color{stewartcyan}\rule{1.15ex}{1.15ex}}\hspace{0.6em}\textbf{\textsf{\large Đạo hàm cấp hai của hàm ẩn}}
    \vspace{0.3em}

    Ví dụ sau đây chỉ ra cách tìm đạo hàm cấp hai của một hàm số được xác định ẩn.

    \vspace{0.6em}
    \noindent\textbf{\textsf{\color{stewartcyan}VÍ DỤ 4}}\hspace{1em}Tìm $y''$ nếu $x^4 + y^4 = 16$.

    \vspace{0.3em}
    \noindent\textbf{\textsf{\color{stewartcyan}LỜI GIẢI}}\hspace{1em}Lấy đạo hàm ẩn hai vế của phương trình theo $x$, ta được
    \begin{equation*}
    4x^3 + 4y^3y' = 0
    \end{equation*}
    Giải tìm $y'$ cho ta
    \vspace{0.2em}
    \begin{center}
    \noindent\makebox[0pt][l]{\fcolorbox{stewartred}{white}{\bfseries\sffamily\color{stewartred}\footnotesize\ 3\ }}%
    \hfill$\displaystyle y' = -\frac{x^3}{y^3}$\hfill\mbox{}
    \end{center}
\end{minipage}

\end{document}